\subsection{Solving Equations} 
\begin{definition}[Solution of an Equation]
A \term{solution} of an equation is a value (or values) of the variable (or variables) that make the equation \makeblank[1in].
\end{definition}
\begin{Example}
Confirm that $x=2$ is a solution to the equation $3x^2-5x-2=0$.
\begin{cpic}{}{2}
\end{cpic}
\end{Example}
\begin{definition}[Solve an Equation]
To \term{solve} an equation means to \emph{find} all the solutions to an equation, and also to \emph{demonstrate} that we have found all the solutions to the equation. To show that we have all the solutions, we must ``\makeblank.''
\end{definition}
Sometimes, we can solve an equation by \emph{inspection}: we look at the equation, and know the solution.
\begin{Example} Solve $x+4=10$.
\begin{itemize}
\item We know that $\makeblanksmall+4=10$, so the only solution must be $x=\makeblanksmall$.
\end{itemize}
\end{Example}
Otherwise, we need a process. We will expand on this process throughout the course.
\begin{displaybox}[Solve Any Equation]
To solve any equation, we might make the following moves.
\begin{enumerate}
\item Rewrite either side of the equation in different forms.
\item Use properties of equality to isolate parts of an expression.
\item ???
\item ???
\end{enumerate}
\end{displaybox}
\begin{Example}
Using the first two moves, solve the equation
\[
3(5x-7)=9
\]
in two different ways.
\begin{lpic}{}{4}
\twocolleft{-1}{4}{Method 1:}{Method 2:}
\end{lpic}
\end{Example}